
%-------------------------------------------------%
% Apresentação Markupp
%-------------------------------------------------%
\documentclass[aspectratio=169]{beamer}

\usepackage[english,brazil]{babel}

\usepackage{tikz}
\usetikzlibrary{positioning, arrows.meta, shapes.geometric, fit, backgrounds}

\definecolor[named]{novaCor}{HTML}{303A52}
\definecolor{azul-seg}{RGB}{27, 85, 142}

\usetheme[
    foreground=azul-seg,
]{ifsclean}

% -------------------------------------------------%
\title{Markupp}
\subtitle{Apresentação final do MVP, v0.4.0}
\author{
  Renato Freitas \and Nícolas Arthur \and Nícolas Pitz \and Gabriela Riedel
}
\institute{Instituto Federal de Santa Catarina}
\date{Junho de 2026}

% -------------------------------------------------%
\begin{document}

\begin{frame}[plain, noframenumbering]
    \maketitle
\end{frame}

\begin{frame}[plain, noframenumbering]{Licenciamento}
    \licenciamentoLivre
\end{frame}

% -------------------------------------------------%
\section{Identificação do projeto}
% -------------------------------------------------%

\begin{frame}{Identificação do projeto}
    \begin{itemize}
        \item \textbf{Sistema}: Markupp, servidor \textit{self-hosted} de notas em Markdown, v0.4.0.
        \item \textbf{Equipe}: Renato Freitas, Nícolas Arthur, Nícolas Pitz, Gabriela Riedel.
        \item \textbf{Problema}: notas em Markdown ficam presas a plataformas proprietárias, sem controle do acervo nem interface aberta a outros clientes, em especial agentes de IA.
        \item \textbf{Público}: usuários e times que querem hospedar o próprio acervo de notas, com integração nativa a agentes de IA.
    \end{itemize}
\end{frame}

% -------------------------------------------------%
\section{Proposta e arquitetura}
% -------------------------------------------------%

\begin{frame}{Proposta e arquitetura}
    \begin{columns}[T]
        \begin{column}{0.45\textwidth}
            \small
            \begin{itemize}
                \item Servidor central com API REST como única interface, aberta a humanos e agentes de IA.
                \item Plugin do Obsidian como cliente humano de referência.
                \item Persistência local em SQLite; servidor encapsula o armazenamento.
                \item Distribuição como imagem Docker que sobe servidor e banco em conjunto.
            \end{itemize}
        \end{column}
        \begin{column}{0.52\textwidth}
            \centering
            \begin{tikzpicture}[
                node distance=0.7cm and 0.6cm,
                client/.style={draw=azul-seg, thick, rounded corners, minimum width=2.3cm, minimum height=0.8cm, align=center, font=\small},
                core/.style={draw=azul-seg, very thick, fill=azul-seg!10, rounded corners, minimum width=3.4cm, minimum height=1cm, align=center, font=\small\bfseries},
                store/.style={draw=azul-seg, thick, rounded corners, minimum width=2.5cm, minimum height=0.8cm, align=center, font=\small},
                container/.style={draw=azul-seg!60, thick, dashed, rounded corners, inner sep=0.4cm},
                arrow/.style={-Stealth, thick, azul-seg}
            ]
                \node[client] (obs) {Plugin Obsidian};
                \node[client, right=1.0cm of obs] (ai) {Agente de IA};
                \node[core, below=of $(obs)!0.5!(ai)$] (api) {API REST (Go)};
                \node[store, below=of api] (db) {SQLite};

                \draw[arrow, <->] (obs) -- (api);
                \draw[arrow, <->] (ai) -- (api);
                \draw[arrow, <->] (api) -- (db);

                \begin{scope}[on background layer]
                    \node[container, fit=(api) (db), label={[font=\scriptsize\itshape, azul-seg]below:Imagem Docker}] {};
                \end{scope}
            \end{tikzpicture}
        \end{column}
    \end{columns}
\end{frame}

% -------------------------------------------------%
\section{Resumo do MVP}
% -------------------------------------------------%

\begin{frame}{Resumo do MVP: planejado, concluído e fora do escopo}
    \begin{columns}[T]
        \begin{column}{0.48\textwidth}
            \textbf{MVP planejado e concluído}
            \begin{itemize}\small
                \item CRUD de notas no servidor central.
                \item Plugin do Obsidian como cliente humano.
                \item API REST como interface única.
                \item Imagem Docker pública e contrato OpenAPI.
            \end{itemize}
        \end{column}
        \begin{column}{0.48\textwidth}
            \textbf{Fora do escopo do MVP}
            \begin{itemize}\small
                \item Autenticação, TLS e \textit{rate limit} (ADR-0005).
                \item Busca léxica e semântica.
                \item Organização hierárquica por pastas.
                \item Versionamento e histórico de notas.
                \item Integração com agentes de IA via MCP.
                \item Colaboração em tempo real.
            \end{itemize}
        \end{column}
    \end{columns}
\end{frame}

\begin{frame}{Resumo do MVP: mudanças de escopo}
    \begin{itemize}
        \item \textbf{Cliente de referência}: o escopo inicial considerava uma interface web como cliente; a decisão foi adotar um \textbf{plugin do Obsidian} como cliente humano único (ADR-0004), aproveitando editor e navegação já estabelecidos.
        \item \textbf{Identidade do servidor}: a pasta \texttt{backend} foi renomeada para \texttt{markupp} (ADR-0009), refletindo que o servidor é o produto central, não acessório de um cliente.
        \item \textbf{Sincronização do plugin}: comandos avulsos (\textit{upload}, \textit{download}, \textit{import}, \textit{sync}) consolidados em uma \emph{Source Control View} estilo \textit{git} (ADR-0011), com as operações \textit{fetch}, \textit{pull}, \textit{push} e \textit{sync}.
    \end{itemize}
\end{frame}

% -------------------------------------------------%
\section{Demonstração}
% -------------------------------------------------%

\begin{frame}{Demonstração: ambiente e cenário}
    \begin{itemize}
        \item \textbf{Release demonstrada}: v0.4.0.
        \item \textbf{Imagem pública}: \texttt{riedelgab/ifsces2:latest} no Docker Hub.
        \item \textbf{Subida do servidor}: \\
        \quad{\small\texttt{docker run -d -p 8080:8080 -v markupp\_data:/data riedelgab/ifsces2:latest}}
        \item \textbf{Plugin}: instalado de \texttt{markupp-plugin-v0.4.0.zip} no \emph{vault} do Obsidian.
        \item \textbf{Cenário}: servidor recém-iniciado, banco vazio; \emph{vault} com uma nota de exemplo.
    \end{itemize}
\end{frame}

\begin{frame}{Demonstração: fluxo principal de uso}
    \begin{columns}[T]
        \begin{column}{0.42\textwidth}
            \small
            \textbf{Funcionalidade demonstrada}\\
            Criação, sincronização, edição e exclusão de notas, ponta a ponta entre Obsidian e servidor.

            \vspace{0.8em}
            \textbf{Fluxo (criação)}
            \begin{enumerate}\small
                \item Usuário cria nota no Obsidian.
                \item Plugin envia \texttt{POST /notes}.
                \item API grava em SQLite.
                \item Servidor responde \texttt{201 Created}; plugin atualiza a interface.
            \end{enumerate}
        \end{column}
        \begin{column}{0.56\textwidth}
            \centering
            \begin{tikzpicture}[
                actor/.style={draw=azul-seg, thick, rounded corners, fill=azul-seg!10, minimum width=1.6cm, minimum height=0.5cm, align=center, font=\scriptsize\bfseries},
                msg/.style={-Stealth, thick, azul-seg},
                msglabel/.style={azul-seg, font=\tiny, fill=white, inner sep=1pt}
            ]
                \node[actor] (u) at (0, 0) {Usuário};
                \node[actor] (p) at (2.2, 0) {Plugin};
                \node[actor] (a) at (4.4, 0) {API (Go)};
                \node[actor] (d) at (6.6, 0) {SQLite};

                \draw[dashed, azul-seg!40] (u) -- (0, -4.6);
                \draw[dashed, azul-seg!40] (p) -- (2.2, -4.6);
                \draw[dashed, azul-seg!40] (a) -- (4.4, -4.6);
                \draw[dashed, azul-seg!40] (d) -- (6.6, -4.6);

                \draw[msg] (0, -0.7) -- (2.2, -0.7) node[msglabel, midway, above] {cria nota};
                \draw[msg] (2.2, -1.4) -- (4.4, -1.4) node[msglabel, midway, above] {\texttt{POST /notes}};
                \draw[msg] (4.4, -2.1) -- (6.6, -2.1) node[msglabel, midway, above] {\texttt{INSERT}};
                \draw[msg] (6.6, -2.8) -- (4.4, -2.8) node[msglabel, midway, above] {ok};
                \draw[msg] (4.4, -3.5) -- (2.2, -3.5) node[msglabel, midway, above] {\texttt{201 Created}};
                \draw[msg] (2.2, -4.2) -- (0, -4.2) node[msglabel, midway, above] {atualiza UI};
            \end{tikzpicture}
        \end{column}
    \end{columns}
\end{frame}

% -------------------------------------------------%
\section{Situação final}
% -------------------------------------------------%

\begin{frame}{Situação final: o que está funcionando}
    Release \textbf{v0.4.0} consolidada em \texttt{main}.
    \begin{itemize}
        \item CRUD completo de notas via API REST.
        \item Sincronização bidirecional entre plugin do Obsidian e servidor.
        \item Distribuição via imagem Docker pública no Docker Hub.
        \item Plugin do Obsidian distribuído como artefato versionado na release.
        \item Suíte de testes com cobertura acima de 80\% no servidor.
        \item Contrato da API descrito em OpenAPI.
    \end{itemize}
\end{frame}

\begin{frame}{Situação final: limitações e pendências}
    \begin{columns}[T]
        \begin{column}{0.50\textwidth}
            \textbf{Limitações conhecidas}
            \begin{itemize}
                \item Sem autenticação, TLS ou \textit{rate limit}, por decisão (ADR-0005); por consequência, o uso é previsto em rede interna.
                \item Sem upload de anexos binários (imagens e outros arquivos).
                \item Tamanho máximo de nota: 50\,MB.
            \end{itemize}
        \end{column}
        \begin{column}{0.46\textwidth}
            \textbf{Bugs e pendências}
            \begin{itemize}
                \item 17 de 17 \emph{issues} do MVP fechadas.
                \item Nenhum \emph{bug} identificado e não corrigido até o momento.
            \end{itemize}
        \end{column}
    \end{columns}
\end{frame}

\begin{frame}{Situação final: próximos passos}
    Caso o projeto continue, as direções naturais são:
    \begin{itemize}
        \item Autenticação e TLS para exposição pública do servidor.
        \item Upload de anexos binários (imagens e outros arquivos).
        \item Busca léxica em títulos e conteúdo; indexação vetorial e busca semântica.
        \item Organização hierárquica por pastas e projetos.
        \item Versionamento e histórico de notas.
        \item Integração com agentes de IA via MCP.
    \end{itemize}
\end{frame}

\end{document}
